\section{Introduction}
\label{sec:intro}

\subsection{Motivation}
Learning predictive models from electronic health records (EHRs) is constrained by privacy law and institutional policy~\citep{ghassemi2021practical,rieke2020future}. FL enables collaborative model improvement without sharing raw records~\citep{mcmahan2017communication,kairouz2021advances}. Medical FL case studies include imaging~\citep{sheller2020federated} and growing interest in structured clinical data.

\subsection{Problem setting}
We study \textbf{treatment feasibility} prediction from early ICU labs and vitals on MIMIC-IV~\citep{johnson2023mimic}. When physical multi-site data are unavailable, \textbf{disease-stratified nodes} provide a pragmatic simulation of non-independent, non-identically distributed (non-IID) partitions: sepsis, cardiac, respiratory, renal, and diabetes cohorts differ in lab distributions and outcome rates.

\subsection{Contributions}
\begin{itemize}
  \item A fully specified \textbf{five-node federated few-shot} protocol with fixed splits and open exports (\texttt{exp3\_results.csv}).
  \item Empirical \textbf{0.005 AUROC gap} between centralized and federated few-shot training in our harness.
  \item Critical discussion of \textbf{simulation limits} (vs.\ real hospitals) and \textbf{ceiling effects} when comparing to centralized tabular SOTA.
\end{itemize}

\subsection{Organization}
Section~\ref{sec:related} reviews FL and clinical applications. Section~\ref{sec:prelim} formalizes FedAvg. Section~\ref{sec:methods} describes the few-shot client objective and server aggregation. Sections~\ref{sec:exp}--\ref{sec:results} give experiments and results. Sections~\ref{sec:discuss}--\ref{sec:conc} discuss limitations and conclude.
